\chapter{The Quantum Harmonic Oscillator and Its Ladder of States}

\section{The Universal Model Near Stable Equilibrium}

This is the last lecture in the present sequence, though not the last word on
quantum mechanics. Angular momentum is the other natural closing subject, but
the harmonic oscillator fills the available time and angular momentum will
have to wait for the later discussion of quantum fields.

The oscillator deserves that time because it is not merely one special
spring problem. Whenever a system has a stable equilibrium and is displaced
slightly from it, the first nontrivial approximation to its motion is usually
harmonic. A weight suspended from a spring is the obvious model. Molecular
vibrations, sound modes, electromagnetic modes, and small disturbances of
many more complicated systems share the same mathematics.

Consider the hanging mass. Gravity pulls downward with an approximately
constant force, while the spring force grows with extension. At one position
the two balance. If \(x\) measures displacement from that equilibrium rather
than from the spring's unstretched length, gravity disappears from the
equation of small motion and the restoring force is
\begin{equation}
  F(x)=-kx.
  \label{eq:l10-Hooke}
\end{equation}
The minus sign says that either displacement produces a force back toward
\(x=0\).

\begin{figure}[tbp]
  \centering
  \includegraphics[width=0.64\textwidth]{lecture_10_figure_02.png}
  \caption{A suspended mass and spring provide the concrete oscillator model;
  the coordinate is measured from the force-balanced equilibrium.}
  \lectureframe{00:02:02}
\end{figure}

The same conclusion follows for a general potential \(V(q)\) with a
nondegenerate minimum at \(q=q_0\). Put \(x=q-q_0\). Taylor expansion gives
\begin{equation}
  V(q_0+x)
  =
  V(q_0)
  +\underbrace{V'(q_0)}_{0}x
  +\frac12 V''(q_0)x^2
  +O(x^3).
  \label{eq:l10-Taylor-potential}
\end{equation}
The constant does not affect the force, the linear term vanishes at the
minimum, and the positive quadratic coefficient is the spring constant:
\(k=V''(q_0)>0\). Harmonic motion is therefore the local language of stable
equilibrium.

\section{Classical Preparation}

Before quantizing, let us put the classical system into a convenient form.
For mass \(m\) and angular frequency
\(\omega=\sqrt{k/m}\),
\begin{equation}
  L
  =
  \frac12m\dot q^{\,2}
  -
  \frac12m\omega^2q^2.
  \label{eq:l10-dimensional-L}
\end{equation}
The rescaling \(x=\sqrt m\,q\) absorbs the mass:
\begin{equation}
  L
  =
  \frac12\dot x^{\,2}
  -
  \frac12\omega^2x^2.
  \label{eq:l10-unit-mass-L}
\end{equation}
We shall use this unit-mass coordinate through most of the lecture and
restore \(m\) whenever its physical role matters.

\begin{figure}[tbp]
  \centering
  \includegraphics[width=0.70\textwidth]{lecture_10_figure_03.png}
  \caption{The spring model and its unit-mass Lagrangian, with a quadratic
  potential characterized by \(\omega^2\).}
  \lectureframe{00:10:57}
\end{figure}

The Euler--Lagrange equation gives
\begin{equation}
  \frac{d}{dt}\frac{\partial L}{\partial\dot x}
  =
  \frac{\partial L}{\partial x},
\end{equation}
or
\begin{equation}
  \boxed{\ddot x=-\omega^2x.}
  \label{eq:l10-classical-eom}
\end{equation}
Its general real solution is
\begin{equation}
  x(t)=A\cos\omega t+B\sin\omega t.
  \label{eq:l10-classical-solution}
\end{equation}
Twice differentiating either sinusoid returns the function multiplied by
\(-\omega^2\), so the frequency in the potential is exactly the frequency
of the motion.

\begin{figure}[tbp]
  \centering
  \includegraphics[width=0.72\textwidth]{lecture_10_figure_04.png}
  \caption{The restoring-force equation and its sinusoidal classical
  solution.}
  \lectureframe{00:12:52}
\end{figure}

The canonical momentum is
\begin{equation}
  p=\frac{\partial L}{\partial\dot x}=\dot x,
  \label{eq:l10-canonical-p}
\end{equation}
and the Legendre transform produces
\begin{align}
  H
  &=p\dot x-L\\
  &=\frac12p^2+\frac12\omega^2x^2.
  \label{eq:l10-classical-H}
\end{align}
The Lagrangian is kinetic minus potential energy; the Hamiltonian is their
sum.

\begin{figure}[tbp]
  \centering
  \includegraphics[width=0.70\textwidth]{lecture_10_figure_05.png}
  \caption{The Legendre transform rewrites the oscillator energy in the
  canonical variables \(x\) and \(p\).}
  \lectureframe{00:17:18}
\end{figure}

\section{Quantization and the Two Jobs of the Hamiltonian}

The quantum state space consists of square-integrable wave functions
\(\psi(x)\). At a specified time,
\begin{equation}
  |\psi(x)|^2\,dx
\end{equation}
is the probability of finding the oscillator coordinate in an interval
\(dx\), and normalization requires
\begin{equation}
  \int_{-\infty}^{\infty}|\psi(x)|^2\,dx=1.
  \label{eq:l10-normalization}
\end{equation}

In the position representation,
\begin{equation}
  (\hat x\psi)(x)=x\psi(x),
  \qquad
  (\hat p\psi)(x)=-i\hbar\frac{\partial\psi}{\partial x}.
  \label{eq:l10-xp-representation}
\end{equation}
Writing \(\hat x\ket\psi\) and writing multiplication by \(x\) are not two
different operations; the second is the position-space representation of
the first. Similarly, \(\hat p^{\,2}\) means applying the momentum operator
twice.

Promoting the classical variables in
Equation~\eqref{eq:l10-classical-H} to operators gives
\begin{equation}
  \hat H
  =
  \frac12\hat p^{\,2}
  +
  \frac12\omega^2\hat x^{\,2}.
  \label{eq:l10-quantum-H}
\end{equation}
Its action on a wave function is
\begin{equation}
  (\hat H\psi)(x)
  =
  -\frac{\hbar^2}{2}
  \frac{\partial^2\psi}{\partial x^2}
  +
  \frac12\omega^2x^2\psi(x).
  \label{eq:l10-H-action}
\end{equation}

\begin{figure}[tbp]
  \centering
  \includegraphics[width=0.76\textwidth]{lecture_10_figure_06.png}
  \caption{The classical Hamiltonian and its action as a differential
  operator on the wave function.}
  \lectureframe{00:22:35}
\end{figure}

The Hamiltonian has two closely related jobs. First, it generates time
evolution:
\begin{equation}
  i\hbar\frac{\partial\psi(x,t)}{\partial t}
  =
  \left[
    -\frac{\hbar^2}{2}\frac{\partial^2}{\partial x^2}
    +\frac12\omega^2x^2
  \right]\psi(x,t).
  \label{eq:l10-time-Schrodinger}
\end{equation}
Even if \(\psi(x,0)\) happens to be real, the factor \(i\) generally produces
an imaginary component immediately. A nonstationary wave function is
therefore intrinsically complex.

\begin{figure}[tbp]
  \centering
  \includegraphics[width=0.76\textwidth]{lecture_10_figure_07.png}
  \caption{The time-dependent Schr\"odinger equation for the oscillator.}
  \lectureframe{00:24:24}
\end{figure}

\begin{classroomqa}
\audiencequestion{Why are the derivatives partial derivatives rather than
ordinary derivatives?
\lecturetimestamp{00:21:24--00:24:07}}
\lecturerresponse{The wave function depends independently on \(x\) and
\(t\). The derivative \(\partial_t\psi\) changes time while holding position
fixed, whereas \(\partial_x\psi\) changes position while holding time fixed.
Time labels evolution; it does not have the same operator status as position
in this formulation.}
\end{classroomqa}

Given an initial wave function, one could integrate
Equation~\eqref{eq:l10-time-Schrodinger} forward in small time steps. The
second job of \(\hat H\), however, gives a more revealing route. Solve the
stationary eigenvalue problem
\begin{equation}
  \hat H\psi_E=E\psi_E,
  \label{eq:l10-eigen-abstract}
\end{equation}
or explicitly,
\begin{equation}
  -\frac{\hbar^2}{2}\frac{d^2\psi_E}{dx^2}
  +
  \frac12\omega^2x^2\psi_E
  =
  E\psi_E.
  \label{eq:l10-stationary-Schrodinger}
\end{equation}

\begin{figure}[tbp]
  \centering
  \includegraphics[width=0.76\textwidth]{lecture_10_figure_08.png}
  \caption{The stationary Schr\"odinger equation simultaneously determines
  the oscillator eigenfunctions and their allowed energies.}
  \lectureframe{00:29:24}
\end{figure}

Once the eigenfunctions are known, arbitrary initial data may be expanded as
\begin{equation}
  \psi(x,0)=\sum_n c_n\psi_n(x),
\end{equation}
and its time evolution is immediate:
\begin{equation}
  \psi(x,t)
  =
  \sum_n c_n e^{-iE_nt/\hbar}\psi_n(x).
  \label{eq:l10-spectral-evolution}
\end{equation}

\section{Why the Bound-State Energies Are Discrete}

Equation~\eqref{eq:l10-stationary-Schrodinger} is second order, so for a
fixed \(E\) it has two linearly independent local solutions. Far from the
origin, the quadratic potential dominates. The two asymptotic behaviors are
of the form
\begin{equation}
  \psi(x)\sim
  \exp\left(\pm\frac{\omega x^2}{2\hbar}\right).
  \label{eq:l10-asymptotic}
\end{equation}
One decays and one grows. For a generic trial energy, a solution adjusted to
decay at \(+\infty\) contains a growing component at \(-\infty\), or vice
versa. Only selected energies allow the growing component to vanish at both
ends.

\begin{figure}[tbp]
  \centering
  \includegraphics[width=0.70\textwidth]{lecture_10_figure_09.png}
  \caption{Normalizability rejects the exponentially growing asymptotic
  branch and selects isolated energies that decay at both ends.}
  \lectureframe{00:34:20}
\end{figure}

This is the origin of the discrete bound-state spectrum. The differential
equation exists for arbitrary \(E\), but a physical state must belong to the
Hilbert space. A function that grows exponentially cannot satisfy
Equation~\eqref{eq:l10-normalization}; in the oscillator potential it also
has divergent potential-energy expectation.

There is a second way to anticipate the lowest energy. Classically one
minimizes
\begin{equation}
  H=\frac12p^2+\frac12\omega^2x^2
\end{equation}
by setting \(x=p=0\). Quantum mechanically,
\begin{equation}
  \Delta x\,\Delta p\geq\frac{\hbar}{2},
  \label{eq:l10-uncertainty}
\end{equation}
so no normalized state can make both contributions vanish. The ground-state
energy must be positive.

\section{Finding the Gaussian Ground State}

The asymptotic form suggests a Gaussian:
\begin{equation}
  \psi_0(x)=C\exp\left(-\frac{\omega x^2}{2\hbar}\right),
  \label{eq:l10-ground-ansatz}
\end{equation}
where \(\omega>0\) denotes the positive oscillator frequency. Its first two
derivatives are
\begin{align}
  \psi_0'
  &=-\frac{\omega x}{\hbar}\psi_0,\\
  \psi_0''
  &=
  \left(
    \frac{\omega^2x^2}{\hbar^2}
    -\frac{\omega}{\hbar}
  \right)\psi_0.
  \label{eq:l10-ground-derivatives}
\end{align}

\begin{figure}[tbp]
  \centering
  \includegraphics[width=0.66\textwidth]{lecture_10_figure_10.png}
  \caption{The Gaussian trial state is concentrated near the stable
  equilibrium and decays in both directions.}
  \lectureframe{00:38:18}
\end{figure}

\begin{figure}[tbp]
  \centering
  \includegraphics[width=0.74\textwidth]{lecture_10_figure_11.png}
  \caption{Differentiating the Gaussian twice produces a constant term and a
  quadratic term.}
  \lectureframe{00:41:05}
\end{figure}

Applying the Hamiltonian,
\begin{align}
  \hat H\psi_0
  &=
  -\frac{\hbar^2}{2}
  \left(
    \frac{\omega^2x^2}{\hbar^2}
    -\frac{\omega}{\hbar}
  \right)\psi_0
  +
  \frac12\omega^2x^2\psi_0\\
  &=
  \frac12\hbar\omega\,\psi_0.
  \label{eq:l10-ground-eigen}
\end{align}
The two terms proportional to \(x^2\) cancel, leaving
\begin{equation}
  \boxed{E_0=\frac12\hbar\omega.}
  \label{eq:l10-zero-point}
\end{equation}

\begin{figure}[tbp]
  \centering
  \includegraphics[width=0.76\textwidth]{lecture_10_figure_12.png}
  \caption{The Gaussian eigenfunction calculation leaves the zero-point
  energy after the quadratic terms cancel.}
  \lectureframe{00:44:34}
\end{figure}

Normalization fixes
\begin{equation}
  C=\left(\frac{\omega}{\pi\hbar}\right)^{1/4},
  \label{eq:l10-ground-normalization}
\end{equation}
up to an irrelevant overall phase. For a physical coordinate \(q\) with
mass \(m\), the replacements \(x=\sqrt m\,q\) give
\begin{equation}
  \psi_0(q)
  =
  \left(\frac{m\omega}{\pi\hbar}\right)^{1/4}
  \exp\left(-\frac{m\omega q^2}{2\hbar}\right).
  \label{eq:l10-dimensional-ground}
\end{equation}

\begin{classroomqa}
\audiencequestion{Does the uncertainty principle mean that the oscillator
can never have zero energy?
\lecturetimestamp{00:45:23--00:46:13}}
\lecturerresponse{Yes. A zero-energy state would require both positive terms
\(\langle p^2\rangle/2\) and
\(\omega^2\langle x^2\rangle/2\) to vanish. That would make both position
and momentum sharp at zero, contradicting
Equation~\eqref{eq:l10-uncertainty}. The lowest allowed value is
\(\hbar\omega/2\).}
\end{classroomqa}

\begin{classroomqa}
\audiencequestion{Why not take \(\psi(x)=1\) as a zero-energy state?
\lecturetimestamp{00:46:13--00:47:49}}
\lecturerresponse{A constant wave function is a formal zero-momentum
eigenfunction, so the kinetic term vanishes, but it is neither normalizable
on the line nor localized at \(x=0\). The potential term produces
\(\tfrac12\omega^2x^2\psi\), which is not zero and is not proportional to
the constant state. It therefore does not solve the oscillator eigenvalue
equation.}
\end{classroomqa}

\begin{classroomqa}
\audiencequestion{For \(E=0\), can some linear combination of the two
second-order solutions remain bounded at both ends?
\lecturetimestamp{00:48:02--00:50:27}}
\lecturerresponse{No. Initial data can be tuned so the solution decays in one
direction, but at \(E=0\) it grows in the other. The exceptional tuning that
removes the growing branch at both ends occurs only at the discrete
eigenvalues, the lowest of which is \(\hbar\omega/2\).}
\end{classroomqa}

\begin{classroomqa}
\audiencequestion{Why must the \(\omega\) in the Gaussian exponent be the
positive root?
\lecturetimestamp{00:50:33--00:51:15}}
\lecturerresponse{The classical equation contains only \(\omega^2\), so its
sign is irrelevant there. In the wave function,
\(\exp[-\omega x^2/(2\hbar)]\) decays only for \(\omega>0\). Choosing the
positive frequency is therefore the normalizable branch.}
\end{classroomqa}

Restoring \(\hbar\) is essential. Frequency has units of inverse time, while
energy requires \(\hbar\omega\). For a macroscopic spring the spacing is
usually far too small to notice, and the zero-point energy disappears
relative to ordinary classical scales as \(\hbar\to0\).

\section{Factorizing the Hamiltonian}

The Gaussian gives one eigenstate. To find all of them, factor the sum of
squares. Define the unnormalized operators
\begin{equation}
  \hat b_-=\hat p-i\omega\hat x,
  \qquad
  \hat b_+=\hat p+i\omega\hat x
  =\hat b_-^\dagger.
  \label{eq:l10-b-operators}
\end{equation}
Classically,
\((p+i\omega x)(p-i\omega x)=p^2+\omega^2x^2\).
For noncommuting quantum operators, the chosen order contributes an
additional term:
\begin{align}
  \hat b_+\hat b_-
  &=
  (\hat p+i\omega\hat x)
  (\hat p-i\omega\hat x)\\
  &=
  \hat p^{\,2}+\omega^2\hat x^{\,2}
  +i\omega[\hat x,\hat p]\\
  &=
  \hat p^{\,2}+\omega^2\hat x^{\,2}
  -\hbar\omega.
  \label{eq:l10-b-product}
\end{align}
Hence
\begin{equation}
  \boxed{
  \hat H
  =
  \frac12\hat b_+\hat b_-
  +\frac12\hbar\omega.}
  \label{eq:l10-H-b}
\end{equation}

\begin{figure}[tbp]
  \centering
  \includegraphics[width=0.74\textwidth]{lecture_10_figure_13.png}
  \caption{The classical-looking factorization must retain operator order in
  quantum mechanics.}
  \lectureframe{00:55:24}
\end{figure}

\begin{figure}[tbp]
  \centering
  \includegraphics[width=0.76\textwidth]{lecture_10_figure_14.png}
  \caption{The noncommuting cross terms generate the additive zero-point
  contribution in the factorized Hamiltonian.}
  \lectureframe{00:59:31}
\end{figure}

The first-order operator \(\hat b_-\) annihilates the Gaussian:
\begin{align}
  \hat b_-\psi_0
  &=
  \left(
    -i\hbar\frac{d}{dx}
    -i\omega x
  \right)\psi_0\\
  &=
  \left(i\omega x-i\omega x\right)\psi_0=0.
  \label{eq:l10-b-annihilates}
\end{align}
Therefore the product \(\hat b_+\hat b_-\) contributes zero on the ground
state, and Equation~\eqref{eq:l10-H-b} immediately returns
\(E_0=\hbar\omega/2\).

\begin{figure}[tbp]
  \centering
  \includegraphics[width=0.76\textwidth]{lecture_10_figure_15.png}
  \caption{The lower first-order factor annihilates the Gaussian ground
  state.}
  \lectureframe{01:03:56}
\end{figure}

Normalize the operators:
\begin{equation}
  \hat a=\frac{\hat b_-}{\sqrt{2\hbar\omega}},
  \qquad
  \hat a^\dagger
  =\frac{\hat b_+}{\sqrt{2\hbar\omega}}.
  \label{eq:l10-a-definitions}
\end{equation}
Since
\begin{equation}
  [\hat b_-,\hat b_+]=2\hbar\omega,
\end{equation}
we obtain
\begin{equation}
  \boxed{[\hat a,\hat a^\dagger]=1.}
  \label{eq:l10-a-commutator}
\end{equation}
Equivalently,
\([\hat a^\dagger,\hat a]=-1\). This resolves the sign uncertainty that
arises if the order of the two operators is changed while they are being
named.

\begin{figure}[tbp]
  \centering
  \includegraphics[width=0.76\textwidth]{lecture_10_figure_16.png}
  \caption{Rescaling the first-order factors produces dimensionless ladder
  operators with a unit commutator.}
  \lectureframe{01:08:22}
\end{figure}

The Hamiltonian is now
\begin{equation}
  \boxed{
  \hat H
  =
  \hbar\omega
  \left(
    \hat a^\dagger\hat a+\frac12
  \right).}
  \label{eq:l10-H-number}
\end{equation}
The operator
\begin{equation}
  \hat N=\hat a^\dagger\hat a
  \label{eq:l10-number-operator}
\end{equation}
is positive:
\begin{equation}
  \bra\psi\hat N\ket\psi
  =
  \|\hat a\ket\psi\|^2\geq0.
  \label{eq:l10-N-positive}
\end{equation}
This factorization proves, rather than merely guesses, that no state lies
below \(\hbar\omega/2\).

\begin{figure}[tbp]
  \centering
  \includegraphics[width=0.74\textwidth]{lecture_10_figure_17.png}
  \caption{The oscillator Hamiltonian reduced to the number operator plus
  the zero-point half.}
  \lectureframe{01:09:37}
\end{figure}

\begin{classroomqa}
\audiencequestion{Which commutator has sign \(+1\), and which operator really
annihilates the ground state?
\lecturetimestamp{01:11:05--01:14:09}}
\lecturerresponse{With the consistent definitions in
Equation~\eqref{eq:l10-a-definitions},
\(\hat a\propto\hat p-i\omega\hat x\) annihilates the Gaussian,
\([\hat a,\hat a^\dagger]=1\), and therefore
\([\hat a^\dagger,\hat a]=-1\). The extended board discussion mixed the
naming order, not the physics. Testing the operator directly on
\(\psi_0\) fixes the convention unambiguously.}
\end{classroomqa}

Writing the ground state as \(\ket0\),
\begin{equation}
  \hat a\ket0=0,
  \qquad
  \hat N\ket0=0,
  \qquad
  \hat H\ket0=\frac12\hbar\omega\ket0.
  \label{eq:l10-ground-algebra}
\end{equation}

\begin{figure}[tbp]
  \centering
  \includegraphics[width=0.66\textwidth]{lecture_10_figure_18.png}
  \caption{The algebraic definition of the ground state:
  \(\hat a\ket0=0\).}
  \lectureframe{01:14:47}
\end{figure}

\section{Climbing the Energy Ladder}

Suppose
\begin{equation}
  \hat N\ket n=n\ket n.
  \label{eq:l10-N-eigen}
\end{equation}
Use the commutator to move \(\hat N\) past the raising operator:
\begin{align}
  \hat N\hat a^\dagger
  &=
  \hat a^\dagger\hat a\hat a^\dagger\\
  &=
  \hat a^\dagger
  \left(\hat a^\dagger\hat a+1\right)\\
  &=
  \hat a^\dagger(\hat N+1).
  \label{eq:l10-raising-identity}
\end{align}
Acting on \(\ket n\),
\begin{equation}
  \hat N(\hat a^\dagger\ket n)
  =
  (n+1)\hat a^\dagger\ket n.
  \label{eq:l10-raising-proof}
\end{equation}
Thus \(\hat a^\dagger\ket n\) is an eigenstate with number eigenvalue
\(n+1\), unless it vanishes.

\begin{figure}[tbp]
  \centering
  \includegraphics[width=0.76\textwidth]{lecture_10_figure_19.png}
  \caption{Commuting the number operator through \(\hat a^\dagger\) raises
  its eigenvalue by one.}
  \lectureframe{01:19:29}
\end{figure}

The first excited state is proportional to
\(\hat a^\dagger\ket0\) and has energy
\begin{equation}
  E_1
  =
  \hbar\omega
  \left(1+\frac12\right)
  =
  \frac32\hbar\omega.
  \label{eq:l10-first-energy}
\end{equation}

\begin{figure}[tbp]
  \centering
  \includegraphics[width=0.74\textwidth]{lecture_10_figure_20.png}
  \caption{One application of the raising operator adds
  \(\hbar\omega\) to the ground-state energy.}
  \lectureframe{01:20:00}
\end{figure}

Repeating the argument generates
\begin{equation}
  \ket n
  =
  \frac{(\hat a^\dagger)^n}{\sqrt{n!}}\ket0,
  \qquad n=0,1,2,\ldots ,
  \label{eq:l10-number-states}
\end{equation}
with the normalized ladder actions
\begin{equation}
  \hat a^\dagger\ket n=\sqrt{n+1}\ket{n+1},
  \qquad
  \hat a\ket n=\sqrt n\ket{n-1}.
  \label{eq:l10-ladder-actions}
\end{equation}
At the bottom, \(\hat a\ket0=0\), so the ladder cannot continue to negative
\(n\).

\begin{figure}[tbp]
  \centering
  \includegraphics[width=0.76\textwidth]{lecture_10_figure_21.png}
  \caption{The raising theorem maps a number eigenstate to the next
  eigenstate.}
  \lectureframe{01:24:52}
\end{figure}

The complete spectrum is
\begin{equation}
  \boxed{
  E_n=\hbar\omega\left(n+\frac12\right),
  \qquad n=0,1,2,\ldots .}
  \label{eq:l10-spectrum}
\end{equation}
Successive levels are separated by the fixed quantum
\(\hbar\omega\):
\begin{equation}
  E_{n+1}-E_n=\hbar\omega.
\end{equation}

\begin{figure}[tbp]
  \centering
  \includegraphics[width=0.76\textwidth]{lecture_10_figure_22.png}
  \caption{The equally spaced oscillator spectrum beginning at the
  zero-point energy.}
  \lectureframe{01:25:38}
\end{figure}

\begin{classroomqa}
\audiencequestion{Are there other Hamiltonians with an equally spaced
spectrum?
\lecturetimestamp{01:27:35--01:28:54}}
\lecturerresponse{One can construct them. For example, on a particle-on-a-ring
Hilbert space, an operator proportional to angular momentum has equally
spaced eigenvalues \(m\hbar\) with positive and negative integers \(m\).
Used directly as a Hamiltonian, however, that spectrum is unbounded below
and has no stable ground state. The oscillator is distinguished by its
equally spaced ladder together with a bottom.}
\end{classroomqa}

\section{What the Eigenfunctions Mean}

\begin{classroomqa}
\audiencequestion{What physical system does a quantum harmonic oscillator
represent? Is it an electron orbiting an atom?
\lecturetimestamp{01:29:02--01:32:31}}
\lecturerresponse{A Coulomb orbit is not itself a harmonic oscillator. The
general setting is any stable configuration examined through small
displacements. A diatomic molecule vibrating about its equilibrium bond
length is a good example. More broadly,
Equation~\eqref{eq:l10-Taylor-potential} shows why every ordinary
nondegenerate minimum is approximately quadratic. A quantum system is
characterized by measurable probability distributions and spectra, not by a
literal photograph of a microscopic trajectory.}
\end{classroomqa}

\begin{figure}[tbp]
  \centering
  \includegraphics[width=0.64\textwidth]{lecture_10_figure_23.png}
  \caption{Near a stable minimum, a general potential is locally parabolic
  and therefore harmonic to leading order.}
  \lectureframe{01:30:41}
\end{figure}

For the ground state,
\begin{equation}
  |\psi_0(x)|^2
  =
  \sqrt{\frac{\omega}{\pi\hbar}}
  \exp\left(-\frac{\omega x^2}{\hbar}\right).
  \label{eq:l10-ground-density}
\end{equation}
Its variance is
\begin{equation}
  (\Delta x)^2=\frac{\hbar}{2\omega}
  \label{eq:l10-ground-width}
\end{equation}
in the unit-mass coordinate. Increasing \(\omega\) stiffens the restoring
force and narrows the spatial distribution.

\begin{figure}[tbp]
  \centering
  \includegraphics[width=0.76\textwidth]{lecture_10_figure_24.png}
  \caption{The Gaussian ground-state probability is more tightly localized
  for a stiffer oscillator.}
  \lectureframe{01:32:38}
\end{figure}

Apply the raising operator once:
\begin{align}
  \psi_1(x)
  &\propto
  \left(
    -\hbar\frac{d}{dx}
    +\omega x
  \right)
  e^{-\omega x^2/(2\hbar)}\\
  &\propto
  x\,e^{-\omega x^2/(2\hbar)}.
  \label{eq:l10-first-wavefunction}
\end{align}
It is odd, changes sign at the origin, and has one node. Multiplying a state
by an overall constant or phase, including a minus sign, does not change its
physical predictions.

\begin{figure}[tbp]
  \centering
  \includegraphics[width=0.70\textwidth]{lecture_10_figure_25.png}
  \caption{The first excited state is an odd Gaussian times \(x\), with one
  node at the origin.}
  \lectureframe{01:37:30}
\end{figure}

Repeated raising produces a polynomial times the same Gaussian:
\begin{equation}
  \psi_n(x)
  =
  \frac{1}{\sqrt{2^n n!}}
  \left(\frac{\omega}{\pi\hbar}\right)^{1/4}
  H_n\left(\sqrt{\frac{\omega}{\hbar}}x\right)
  e^{-\omega x^2/(2\hbar)},
  \label{eq:l10-Hermite-wavefunctions}
\end{equation}
where \(H_n\) is the \(n\)th Hermite polynomial. The parity alternates,
\begin{equation}
  \psi_n(-x)=(-1)^n\psi_n(x),
\end{equation}
and \(\psi_n\) has \(n\) real nodes. Higher states oscillate more rapidly,
carry larger characteristic momentum, and extend farther from the origin,
as expected for a more energetic oscillator.

\begin{figure}[tbp]
  \centering
  \includegraphics[width=0.70\textwidth]{lecture_10_figure_26.png}
  \caption{Successive oscillator states acquire additional nodes while
  retaining Gaussian decay at large distance.}
  \lectureframe{01:39:12}
\end{figure}

\begin{classroomqa}
\audiencequestion{Are the nodes of the Hermite-polynomial wave functions at
real values of \(x\), and is the probability exactly zero there?
\lecturetimestamp{01:41:44--01:42:58}}
\lecturerresponse{Yes. The \(n\)th oscillator eigenfunction has \(n\) simple
real zeros. At each node \(\psi_n(x)=0\), so the position probability density
\(|\psi_n(x)|^2\) also vanishes. The zeros arise from destructive
interference within the stationary wave function.}
\end{classroomqa}

\begin{classroomqa}
\audiencequestion{What is the physical reason for quantized energy rather
than a continuum?
\lecturetimestamp{01:42:59--01:44:35}}
\lecturerresponse{The mathematical and physical statements are the same
boundary condition viewed from two sides. The oscillator is confined:
its potential grows without bound, so a physical eigenfunction must decay at
both infinities. Those two boundary requirements can be met only at discrete
energies. A free particle can escape to infinity and therefore has a
continuous momentum and energy spectrum. More general potentials may have
discrete levels below an escape threshold and a continuum above it.}
\end{classroomqa}

\begin{classroomqa}
\audiencequestion{Why are the exponentially growing solutions excluded
rather than treated as unusual states?
\lecturetimestamp{01:44:38--01:46:24}}
\lecturerresponse{They cannot be normalized to a probability distribution on
the line. Relative to their divergent weight at infinity, every finite
region would have zero probability, and the oscillator potential would give
them divergent energy. They are mathematical solutions of the local
differential equation, but not state vectors in the oscillator Hilbert
space.}
\end{classroomqa}

\section{The Completed Oscillator Picture}

The harmonic oscillator begins with the most ordinary classical motion:
\(\ddot x=-\omega^2x\). Quantization turns its coordinate and momentum into
noncommuting operators. That single change prevents zero energy, selects a
normalizable Gaussian ground state, and leaves a half-quantum of irreducible
energy.

Factorization then converts a differential-equation problem into an algebra
problem:
\begin{equation}
  [\hat a,\hat a^\dagger]=1,
  \qquad
  \hat H=\hbar\omega
  \left(\hat a^\dagger\hat a+\frac12\right).
\end{equation}
The creation operator climbs an equally spaced spectrum, the annihilation
operator descends it, and the ground state terminates the descent. In
position space the same ladder appears as Gaussian wave functions multiplied
by successively higher Hermite polynomials.

Angular momentum has not been covered here; it is deferred rather than
silently omitted. The oscillator itself supplies the bridge to the next
subject, because every normal mode of a quantum field behaves as a quantum
harmonic oscillator, and its raising operator creates one additional quantum
of excitation.
